% ===== sec/01_introduction.tex =====
\section{Introduction}
\label{sec:introduction}

\IEEEPARstart{A}{gent} architectures are rich in diagrams and short
on formal definitions.  An agent is variously described as ``a
system that perceives and acts'', ``a planner that calls tools'', or
``an LLM in a loop''---each admitting neither a decision procedure
nor a complexity bound.  Two agents with radically different
computational power are routinely compared as if they belonged to a
single class.

Koohestani et al.~\cite{koohestani2025automata} recently posed a
corrective question: \emph{Are agents just automata?}  Their answer
is affirmative, establishing a memory-architecture equivalence
(simple-reflex~$\cong$~FA, hierarchical~$\cong$~PDA,
unrestricted~$\cong$~TM).  We accept their answer but ask the
operational follow-up: \emph{how can the Chomsky type of one agent
bound another?}  This is the producer--consumer problem of
Dijkstra~\cite{dijkstra1968cooperating} re-stated in the language of
formal grammars.

\subsection{Three claims}

The paper makes three claims:

\begin{enumerate}
\item \textbf{Agent behaviour can be modelled as a trace language.}
An agent emits a finite string of operation symbols; the collection
of all such strings is its trace language
(Section~\ref{sec:trace}).
\item \textbf{A regular (Type-3) verifier can constrain a more
expressive generator while remaining computationally tractable.}
Intersecting a generator's trace language with a regular verifier
language preserves the generator's Chomsky class, and conformance
checking runs in linear time on the verifier side
(Theorem~\ref{thm:closure}).
\item \textbf{The framework works on a realistic benchmark.}
Instantiated on \textsc{Spaceship Titanic} from \textsc{MLE-bench},
the architecture achieves leaderboard placement $710/2330$ (top
tertile), and a controlled ablation shows the regular verifier
raises mean task score from $35.2\%$ to $57.8\%$ while eliminating
undetected trace faults.
\end{enumerate}

\subsection{Related work}

Koohestani et al.~\cite{koohestani2025automata} classify agents by
memory architecture but do not prescribe how classes interact.
Formal-LLM~\cite{li2024formalllm} uses a context-free grammar to
constrain plan generation; our construction is dual---we verify
emission through a \emph{lower}-class recogniser.
Dijkstra's~\cite{dijkstra1968cooperating} producer--consumer framing
is the operational form of our question: what Chomsky-class
discipline on the consumer renders the producer's output decidable?
MLE-bench~\cite{koohestani2025automata} supplies the empirical
benchmark.

\subsection{Contributions}

\begin{enumerate}
\item A trace-language definition of agent behaviour
(Section~\ref{sec:trace}).
\item A producer--consumer answer: a single idealised generator
intersected with a regular verifier preserves the generator class
and keeps conformance checking linear (Theorem~\ref{thm:closure}).
\item An Aho--Corasick~\cite{aho1975corasick} DFA verifier
(Proposition~\ref{prop:verifier-dfa}).
\item A demonstration that intermediate sub-agent roles are
projections, not primitives (Theorem~\ref{thm:derived}).
\item An empirical instantiation on \textsc{Spaceship Titanic}
achieving $710/2330$ and a controlled ablation isolating the
verifier's causal effect (Sections~\ref{sec:expsetup}--\ref{sec:results}).
\end{enumerate}

\subsection{Roadmap}

Section~\ref{sec:trace} defines trace languages and proves the
closure theorem.  Section~\ref{sec:architecture} describes the
generator--verifier binding and projected roles.
Sections~\ref{sec:expsetup}--\ref{sec:results} report the empirical
setup and outcomes.  Section~\ref{sec:implications} distils design
implications and a reusable framework.
Section~\ref{sec:discussion} discusses limitations and threats to
validity.  Supplementary material is available online for extended
theory and additional experiments.
